\documentclass[12pt]{article}
\usepackage[T1]{fontenc}
\usepackage[utf8]{inputenc}
\usepackage{lmodern}
\usepackage{textcomp}
\usepackage{enumitem}
\usepackage{geometry}
\geometry{margin=1in}
\begin{document}
\begin{center}
Answer Key - Philosophy - Undergraduate Level Assessment 10 \
\small (Answers are ordered exactly as in the question paper.)
\end{center}

\section*{Multiple Choice}
\begin{enumerate}[label=\arabic*.]
\item \textbf{Answer:} C
\\ \textit{Options:} A) lie whenever doing so is necessary to prevent a catastrophe.; B) always do whatever brings about the greatest happiness.; C) always assert what you don't think to be the case.; D) all of the above.
\\ \textit{Note:} The chosen answer illustrates a morally impossible rule because asserting something one does not believe contradicts the fundamental principle of honesty, making it ethically untenable.
\item \textbf{Answer:} B
\\ \textit{Options:} A) China; B) Silla; C) Goguryeo; D) Joseon
\\ \textit{Note:} The correct answer is Silla because this ancient Korean kingdom is historically recognized for its reverence of various deities, including the Holy Mother of Mount Fairy Peach, reflecting the region's rich cultural and religious traditions.
\item \textbf{Answer:} A
\\ \textit{Options:} A) one's own will.; B) one's own emotions.; C) God.; D) Nature.; E) the universe.
\\ \textit{Note:} Kant argues that the moral law is a product of rational will, emphasizing that individuals must act according to principles they can will to be universal, thus grounding morality in one's own autonomy and reason.
\item \textbf{Answer:} B
\\ \textit{Options:} A) accomplices in the dictator's crimes.; B) people who knowingly buy stolen goods.; C) thieves who steal from the poor.; D) merchants who sell harmful products.; E) war profiteers exploiting conflict.
\\ \textit{Note:} Singer argues that just as individuals who knowingly purchase stolen goods are complicit in theft and its moral implications, international corporations that engage with corrupt dictators are complicit in the injustices and exploitation perpetuated by those regimes.
\item \textbf{Answer:} A
\\ \textit{Options:} A) strive for success.; B) pursue knowledge.; C) pursue happiness.; D) focus on material wealth.; E) none of the above.
\\ \textit{Note:} Feinberg argues that actively striving for success fosters personal growth, fulfillment, and a sense of purpose, which are essential components in the pursuit of genuine happiness.
\end{enumerate}

\section*{True / False}
\begin{enumerate}[label=\arabic*.]
\setcounter{enumi}{5}
\item \textbf{Answer:} True
\\ \textit{Options:} A) True; B) False
\\ \textit{Note:} The laudatory personality fallacy is true because it illustrates how an individual's admirable characteristics can lead to the overlooking or excusing of their negative behaviors, thereby failing to assess their actions based on merit alone.
\item \textbf{Answer:} False
\\ \textit{Options:} A) True; B) False
\\ \textit{Note:} Velleman argues that the moral complexities and personal nature of euthanasia decisions should be guided by individual autonomy and ethical considerations rather than being solely determined by a court of law.
\item \textbf{Answer:} True
\\ \textit{Options:} A) True; B) False
\\ \textit{Note:} Anscombe argues that moral obligations arise from human desires and preferences, suggesting that our understanding of what we ought to do is rooted in our subjective inclinations and the values we attach to our actions.
\item \textbf{Answer:} True
\\ \textit{Options:} A) True; B) False
\\ \textit{Note:} The fallacy of reprehensible personality involves dismissing an argument based on the character of the person presenting it, which is a specific application of the broader fallacy of accident where general principles are improperly applied to individual cases.
\item \textbf{Answer:} True
\\ \textit{Options:} A) True; B) False
\\ \textit{Note:} Rabbinical commentaries produced after the Mishnah are known as Midrash because they represent interpretative works that expand on biblical texts through analysis, commentary, and storytelling, forming a critical aspect of Jewish exegesis and tradition.
\end{enumerate}

\section*{Long Form Response}
\begin{enumerate}[label=\arabic*.]
\setcounter{enumi}{10}
\item \textbf{Answer:} In his article "Will Cloning Harm People?", Pence employs Kant's humanity formulation by arguing that cloning can be seen as a means to promote human dignity and autonomy. According to Kant, individuals should be treated as ends in themselves, capable of making choices that enhance their lives. Pence suggests that cloning, when conducted ethically, empowers individuals by providing them with reproductive options and potentially alleviating suffering through medical advancements. This perspective challenges traditional bioethical concerns by framing cloning not as a violation of human dignity, but as a means to respect and enhance it. The implications for bioethics are significant, as they encourage a reevaluation of how emerging technologies are perceived in light of Kantian principles, promoting a more nuanced understanding of moral acceptability in biotechnological advancements.
\\ \textit{Note:} correct because it accurately interprets Kant's humanity formulation as promoting human dignity and autonomy, arguing that cloning can enhance individual choices and alleviate suffering, thereby aligning with ethical considerations in bioethics.
\item \textbf{Answer:} The fallacy of appeal to pride occurs when an argument urges acceptance of a claim based solely on the favorable characteristics or status of the individual making it, rather than on the merits of the claim itself. For example, one might assert that a particular political policy is valid because it was proposed by a respected leader, rather than presenting evidence to support the policy's effectiveness. This undermines rational discourse by shifting the focus from the validity of the argument to the personal qualities of the individual, thereby fostering an environment where emotional appeals take precedence over logical reasoning. Consequently, such arguments can lead to the acceptance of claims that lack empirical support, ultimately weakening the critical analysis necessary for sound decision-making in philosophical discussions.
\\ \textit{Note:} The answer correctly identifies the fallacy of appeal to pride as an argument that prioritizes the speaker's status over the actual merits of the claim, illustrating how this shift detracts from rational discourse by neglecting evidence and logical reasoning.
\end{enumerate}

\vfill
\noindent\textit{End of Answer Key}
\end{document}